\section{A uniform local expansion through changes of arithmetic type}
\label{sec:v41-transition}
The fixed-radius arithmetic formula does not describe a sequence of
radii on which a peripheral eigenvalue moves into the unit disk at the
inverse collision scale. We retain that eigenvalue and evaluate its
Fourier integral. The resulting formula is a finite sum of Gaussian
terms with explicitly specified damping, phase and drift. It is uniform
in the radius and does not require a continuous choice of an exact
resonance group.

Write $\mathsf f_R=X_R+\mu_R$, where
$\mu_R=(0,0,\bar\tau_R,c)$ and $X_R$ is
\eqref{eq:v37-occupation-record}. Thus
$\mathsf f_R=(\kappa_{R,1},\kappa_{R,2},\tau_R,\eta_R)$.
For a target put $y=(k_1,k_2,t,n)$, in this coordinate order, and
$Z=(y-m\mu_R)/\sqrt m$. The probability
$P_{n,m,R}(k,t;J)$ retains the exact return meaning of
Theorem~\ref{thm:v40-arithmetic-return-LLT}.

\subsection{Local spectral maxima and their scalar data}
\begin{lemma}[Continuous analytic scalar data]
\label{lem:v41-scalar-continuity}
Near each point of the joint resonance locus, the simple eigenvalue
$\lambda(R,\theta)$ and its logarithm are continuous in $R$, analytic
in the complex frequency coordinates, and have uniformly bounded
frequency derivatives on a smaller neighborhood. Their frequency
derivatives are continuous in $R$ there. The same assertions hold for
\[
 A(R,\theta)=\ell\bigl(M_{\eta_R}\Pi(R,\theta)(\eta_R\nu)\bigr).
\]
These are scalar continuity statements on local common spaces. No
strong-norm differentiability in the table parameter is asserted.
\end{lemma}
\begin{proof}
Use the local rank-one continuation and complementary contour of
Theorems~\ref{thm:v40-full-peripheral} and
\ref{thm:v40-uniform-resonant-reduction}. Shrink its complex frequency
neighborhood so that the continued eigenvalue is nonzero and its
modulus exceeds the complementary contour radius. The contour estimates
and strong frequency analyticity give common bounds for the projection,
the eigenvalue and all derivatives on a smaller polydisk. The
strong-to-weak spectral stability already used in those theorems gives
continuity of the simple eigenvalue at real frequencies. The common
complex bounds make every convergent parameter sequence a normal family
of holomorphic eigenvalues. Every subsequential limit agrees on a real
open frequency ball, and hence on the polydisk, with the limiting
eigenvalue; uniqueness follows by the one-variable identity theorem
successively in each coordinate. Thus convergence is locally uniform
also in complex frequencies. Cauchy's formula gives continuity of the
frequency derivatives. It is not differentiation in the radius.

We justify the moving endpoint amplitude without asserting strong-norm
continuity of the indicator multiplier. For a fixed integer $N$, put
\[
 F_N(R,\theta)=\int\eta_R(x)\eta_R(T_R^Nx)
     e^{i\theta\cdot\sum_{j<N}\mathsf f_R(T_R^jx)}\,d\nu(x).
\]
This finite pairing is continuous in $R$: outside the finitely many
singularity preimages and section boundaries at the limiting radius,
the finite itinerary and all membership decisions persist. Dominated
convergence applies, also uniformly on compact complex frequency sets
for this fixed $N$. The local splitting gives
\[
 \lambda(R,\theta)^{-N}F_N(R,\theta)
     =A(R,\theta)+O(\rho_1^N),\qquad \rho_1<1,
\]
uniformly on a smaller neighborhood. Therefore $A$ is continuous as a
uniform limit of continuous scalar functions. Its holomorphy and common
bounds follow either from this formula or from the projection pairing.
Cauchy's formula again gives derivative continuity. A finite cover of
the compact joint resonance locus supplies the stated uniform bounds.
\end{proof}

\begin{lemma}[Real maxima of the continued branches]
\label{lem:v41-moving-maxima}
One may choose finitely many resonance charts and frequency-smooth
partition functions $\chi_j$, with compact supports in those charts,
so that the spectral complement decays exponentially on each fixed
roof band and the following data are defined on every parameter
projection of a chart:
\begin{align}
 a_j(R)&\in\R^4,\qquad
 z_j(R)=\lambda_j(R,a_j(R)),\qquad
 \kappa_j(R)=-\log|z_j(R)|\ge0,\notag\\
 \mu_j(R)&=\frac1i\nabla_\theta\log\lambda_j(R,a_j(R))\in\R^4,\notag\\
 H_j(R)&=-\nabla_\theta^2\log\lambda_j(R,a_j(R)).
                                      \label{eq:v41-peak-data}
\end{align}
Here $a_j$ is the unique real critical point of
$\log|\lambda_j|$ in its chart. All these data are continuous, $H_j$ is
complex symmetric, and $\operatorname{Re}H_j\ge aI_4$ for a common
$a>0$. In local real frequency coordinates,
\begin{align}
 |\lambda_j(R,a_j+x)|&\le e^{-\kappa_j-a|x|^2},
                                      \label{eq:v41-peak-bound}\\
 \log\lambda_j(R,a_j+x)
  &=\log z_j+i\mu_j\cdot x-\tfrac12x^{\mathsf T}H_jx+O(|x|^3).
                                      \label{eq:v41-peak-taylor}
\end{align}
If $\kappa_j(R)=0$, then $a_j(R)$ is an actual resonance. At such a
parameter its roof coordinate is zero, $\mu_j(R)=\mu_R$, and
$H_j(R)=\Omega_R$.
\end{lemma}
\begin{proof}
At an actual resonance, Proposition~\ref{prop:v40-resonant-curvature}
gives zero real gradient and negative definite real Hessian for
$\log\lambda$. Lemma~\ref{lem:v41-scalar-continuity} makes the
Hessian uniformly negative on a smaller real ball and for nearby
parameters. The equation for its real gradient has a unique nearby
zero. For completeness this parameter version of the implicit function
argument needs only continuity in $R$: multiply the gradient by the
inverse of its Hessian at the central parameter, and use the resulting
uniform contraction on a smaller closed ball. The fixed point is
continuous in $R$ and remains strictly inside that ball.

The real gradient vanishes there, so the complex gradient is purely
imaginary. Taylor's formula gives \eqref{eq:v41-peak-taylor}; integrating
the negative real Hessian along a segment gives
\eqref{eq:v41-peak-bound}, after reducing $a$. All real-frequency
spectral radii are at most one by Theorem~\ref{thm:v40-full-peripheral},
which proves $\kappa_j\ge0$. Equality identifies a peripheral frequency,
so the last assertions follow from
Theorem~\ref{thm:v40-no-roof-resonance} and
Proposition~\ref{prop:v40-resonant-curvature}.

Choose these charts first around the compact joint resonance locus,
and then cover the remaining compact set by nonresonant charts.
A subordinate partition may have common bounds for each fixed number
of frequency derivatives. Each peak is defined on an open parameter
neighborhood of the compact parameter support of its partition
function, including when that function vanishes at the peak. All actual
resonances have roof coordinate zero. Thus the peak charts may be fixed
in one small roof neighborhood, independently of any later fixed outer
band. Enlarging that band adds only a compact exponentially decaying
complement. No limit is taken with the outer band growing.
\end{proof}

For a complex symmetric matrix $H$ with positive real part define
\begin{equation}\label{eq:v41-complex-gaussian}
 G_H(z)=(2\pi)^{-4}\int_{\R^4}
                 e^{-iv\cdot z-v^{\mathsf T}Hv/2}\,dv,
                    \qquad z\in\R^4.
\end{equation}
This converges absolutely. On the matrix set in the preceding lemma,
$G_H$ and every fixed derivative in $z$ are uniformly bounded. For
real $H>0$ it is the ordinary density $g_H$. Put
\begin{align}
 b_j(R)&=(a_j(R))_3,\qquad
 B_j(R)=\chi_j(R,a_j(R))A_j(R,a_j(R)),\notag\\
 \mathcal K_{m,R}(y)
   &=\frac1c\sum_j e^{-ia_j(R)\cdot y}z_j(R)^m B_j(R)
          G_{H_j(R)}\!\left(\frac{y-m\mu_j(R)}{\sqrt m}\right),\notag\\
 \mathcal L_{m,R}(y)&=\operatorname{Re}\mathcal K_{m,R}(y).
                                      \label{eq:v41-transition-kernel}
\end{align}
Only charts whose parameter supports contain $R$ contribute. Different
real lifts of the torus coordinates give the same expression, since
$k_1,k_2,n$ are integers. The data in this formula are values and the
first two frequency derivatives of finitely many scalar spectral
branches. The formula contains no unevaluated Fourier integral of an
operator power; \eqref{eq:v41-complex-gaussian} is an ordinary Gaussian
integral. The peak charts are fixed once as above.

\subsection{Evaluation of the fixed coefficient}
\begin{theorem}[A uniform transition formula at the exact return index]
\label{thm:v41-uniform-transition}
For every fixed bounded interval $J$ of positive length,
\begin{equation}\label{eq:v41-uniform-interval}
 \sup_{R,n,m,k,t}
 \left|m^2P_{n,m,R}(k,t;J)-|J|\mathcal L_{m,R}(k_1,k_2,t,n)\right|
                         \longrightarrow0\qquad(m\to\infty),
\end{equation}
where $m\ge n\ge1$ and the supremum is over all the other targets.
The kernels are uniformly bounded and
$\sup_{R,y}(-\mathcal L_{m,R}(y))_+\to0$ on these targets.
For a fixed integrable roof test with integrable compactly supported
Fourier transform,
\begin{equation}\label{eq:v41-uniform-test}
 m^2P^q_{n,m,R}(k,t)=\left(\int q\right)\mathcal K_{m,R}(y)+o(1)
\end{equation}
uniformly in the same parameters, including complex $q$. Moreover
$\sup_{R,y}|\operatorname{Im}\mathcal K_{m,R}(y)|\to0$.
At each fixed radius the formula reduces, on central compact sets, to
\[
 \mathcal L_{m,R}(y)
       =c\mathfrak a_R(k,n,m)g_{\Omega_R}(Z)+o_R(1).
\]
Thus the uniform formula retains a branch even when its modulus is
strictly below one at the radius in question.
\end{theorem}
\begin{proof}
Use the fixed-count reduction
\eqref{eq:v40-fixed-spectral-inverse}. On chart $j$ set
$\theta=a_j+x$ and then $x=v/\sqrt m$. Equations
\eqref{eq:v41-peak-bound}--\eqref{eq:v41-peak-taylor} give its
contribution, after multiplication by $m^2$, as
\begin{equation}\label{eq:v41-test-peak-sum}
 \frac1c e^{-ia_j\cdot y}z_j^m B_j\widehat q(-b_j)
            G_{H_j}\!\left(\frac{y-m\mu_j}{\sqrt m}\right)+o(1).
\end{equation}
The error is uniform in $R$ and in all targets. Indeed the target
oscillation has modulus one. The real Gaussian bound dominates the
rescaled integrand. On $|v|\le m^{1/12}$ the cubic error is bounded by
$Cm^{-1/2}|v|^3$ times a fixed Gaussian; outside that ball the Gaussian
tail tends to zero uniformly. Frequency variation of the projector
and partition costs $O(m^{-1/2})$ in the integrated estimate.
Continuity of $\widehat q$ pays its remaining variation. Extend the
smooth partition by zero inside the larger chart; the same estimate
then holds if its value at the peak is zero. Summing the finitely many
charts and the exponentially decaying complement proves
\eqref{eq:v41-test-peak-sum} with a uniform total error.

We may replace $\widehat q(-b_j)$ by $\widehat q(0)=\int q$ in this
sum. Here is the uniform argument. On each compact parameter support,
$\kappa_j=0$ implies $b_j=0$. For any fixed $\delta>0$, continuity and
compactness give a positive minimum of $\kappa_j$ on $|b_j|\ge\delta$,
unless that set is empty. Its contribution is exponentially small.
On $|b_j|<\delta$, uniform continuity of $\widehat q$ gives a bound
which tends to zero with $\delta$. First let $m\to\infty$ and then
$\delta\to0$. This proves \eqref{eq:v41-uniform-test}. All prefactors
and Gaussian integrals are bounded, so $|\mathcal K_{m,R}|\le C$.
Taking one fixed real nonnegative band-limited test of integral one
shows that its imaginary part tends to zero uniformly.

For $\one_J$ use the positive-order upper and lower band-limited
envelopes from Theorem~\ref{thm:two-endpoint-collision-LLT}. Their
integrals are $|J|\pm2\pi/B$. The upper envelope is nonnegative; the
lower need not be. Order is preserved by the original exact-index
source. For fixed $B$, \eqref{eq:v41-uniform-test} bounds the upper
and lower expressions by their masses times $\mathcal L_{m,R}$,
with uniform errors. Their difference and $|\mathcal L_{m,R}|\le C$
give a limsup error at most $C/B$. Let $B\to\infty$ only after the
collision limit. Positivity of the actual probability gives the
assertion about the negative part by taking one interval of length
one. This is a qualitative limit, not a polynomial local-limit rate.

At fixed $R$, charts whose peaks have $\kappa_j>0$ decay exponentially.
Every other peak is an actual resonance, where its drift and curvature
are $\mu_R,\Omega_R$. The partition weights at each such point sum
to one. Proposition~\ref{prop:v40-general-residue} and finite Fourier
orthogonality \eqref{eq:v40-finite-residue-sum} now give the last formula.
In contrast, along $R=R_m$ a factor $e^{-m\kappa_j(R_m)}$ need not
vanish merely because every $\kappa_j(R_m)$ is positive.
\end{proof}

\begin{corollary}[A common, slowly shrinking roof interval]
\label{cor:v41-shrinking-roof}
There is a deterministic sequence $h_m\downarrow0$ such that, for
$J_m=[0,h_m]$,
\[
 \sup_{R,n,k,t}
 \left|\frac{m^2}{h_m}P_{n,m,R}(k,t;J_m)-\mathcal L_{m,R}(y)\right|
                       \longrightarrow0.
\]
No specified polynomial rate of decrease for $h_m$ is asserted.
\end{corollary}
\begin{proof}
Apply \eqref{eq:v41-uniform-interval} to $[0,1/j]$ for each integer
$j\ge1$. Choose increasing integers $M_j$ such that its absolute error
is at most $j^{-2}$ whenever $m\ge M_j$. Set $h_m=1/j$ for
$M_j\le m<M_{j+1}$, increasing $M_j$ if necessary. The normalized
error is at most $1/j$. This diagonal uses only fixed-band statements
at each fixed $j$; it supplies neither a density value nor a high-roof
frequency estimate at a prescribed growing band.
\end{proof}
